\section{Methods}
\label{sec:methods}

\subsection{Survey-correct estimation}
Every descriptive statistic and every model metric is weighted by the final student
weight $W_{\mathrm{FSTUWT}}$. Where a model is fit, weights are routed to the
estimator's own \texttt{sample\_weight} (correcting a subtle bug in which
scaler-wrapped models, logistic regression, SVM, KNN, MLP, silently fell back
to an unweighted fit because the weight was passed to the wrong pipeline step).
Class-conditional statistics rescale the population-scale weights to the effective
sample size before, e.g., a $\chi^2$ test, so significance is not inflated by the
weight magnitudes.

\subsection{Leakage-safe pipeline}
Each fold's pipeline is: (1) an engineered-feature builder that learns component
means, standard deviations, and the home-possessions rank curve on the training fold
and replays them on the test fold; (2) median imputation fit on the training fold;
(3) the estimator. Interaction terms are built from mean-centred components so a
product is not confounded with its main effects. School-composition features are
survey-weighted, \emph{leave-one-out} school means (a student's own value is excluded
from their school mean), and a fold-safe re-computation that recomputes the means
from the training fold only reproduces an identical effect, so the lift is
compositional, not leakage.

\subsection{Evaluation and comparison}
Models are scored with repeated stratified $5\times4$ cross-validation, weighted ROC
and PR AUC, F1-macro, and MCC. Fold-score confidence intervals use the
Nadeau--Bengio corrected-resampled variance, which replaces the naive $1/k$ term with
$1/k + n_{\text{test}}/n_{\text{train}}$ to account for the correlation between folds
that share training data. Model-vs-model differences are tested directly: the
corrected resampled $t$-test on paired per-fold AUCs (cross-validation) and the
DeLong test on shared test-set predictions (out-of-sample). This turns a ranking
table into a claim about which gaps are real.

\subsection{Covariate-shift diagnostic}
To quantify the 2018$\rightarrow$2022 change we train a domain classifier to
distinguish the two cohorts from background features alone, with weighted subsampling,
pooled-median imputation, and missingness indicators; a high AUC certifies that the
input distribution, not just the score level, moved. We validate the diagnostic
on a null (same cohort split at random $\rightarrow$ AUC~$\approx$~0.5).

\subsection{Explanation, calibration, fairness, forecast}
We compute TreeSHAP global and per-student attributions on the school-context model,
and partial-dependence / centred-ICE curves for the top drivers. For deployment we
add out-of-fold isotonic calibration and tune the decision threshold on calibrated
out-of-fold predictions to maximise weighted MCC. The fairness audit reports weighted
demographic parity, equal-opportunity (TPR) and false-positive-rate gaps across
gender, SES quintile, and immigration status, plus a threshold sweep. A random-intercept
logistic model provides the design-appropriate specification and the between-school
intraclass correlation. Finally, a scenario Monte-Carlo forecast for 2026 propagates
each cycle's design-based standard error through a weighted trend fit; with only five
cycles and a structural break we report scenarios and their intervals, not a point
prediction.

\subsection{Robustness and extension analyses}
Five further analyses probe the main claims. (i) \emph{Cross-country ceiling
replication}: the student-only versus school-context comparison is re-run identically in
all nine systems, and each system's between-school ICC is read from its random-intercept
model, to test whether the feature ceiling and its size generalise. (ii) \emph{Modality
ablation}: PISA's computer-based cognitive-process indicators are added to the
school-context model and split into cognitive-process and questionnaire-process blocks,
so an endogenous within-test signal can be separated from data available at screening
time. (iii) \emph{Fairness mitigation}: group-specific decision thresholds are fit on
calibrated out-of-fold scores to equalise the per-quintile false-positive rate
(equalized-odds post-processing), tracing the accuracy--equity frontier. (iv)
\emph{Partial identification}: given the 79\% coverage index, worst-case and monotone
Manski bounds on the population at-risk rate are computed, the monotone bound assuming
the out-of-frame group is no better off than the assessed. (v) \emph{Difference-in-differences}:
a design-based DiD contrasts the earthquake-affected central region with the rest of the
country across 2018 and 2022, with a 2015$\rightarrow$2018 placebo to test parallel
trends; weights, plausible values, and replicate-weight standard errors are carried
throughout.

\subsection{Reproducibility and engineering}
Logic lives in unit-tested modules (177 tests on synthetic fixtures, covering the
weighted statistics, Rubin's rules, the leakage-safe transformer, the BRR+PV variance,
the Manski bounds and group-threshold mitigation, and the Nadeau--Bengio / DeLong
maths); notebooks are narrative layers over them. On
macOS the boosters (Homebrew \texttt{libomp}) and scikit-learn (bundled OpenMP) can
abort a process if their runtimes collide; we isolate each model fit in a fresh
subprocess that imports the boosters before scikit-learn, which resolved an
import-order segfault and returned LightGBM to the comparison.
